\documentclass[11pt]{article}
\usepackage[margin=1in]{geometry}
\usepackage{amsmath,amssymb,amsthm,mathtools}
\usepackage{hyperref}

\newtheorem{theorem}{Theorem}
\newtheorem{lemma}{Lemma}
\newcommand{\E}{\mathbb E}
\newcommand{\R}{\mathbb R}

\title{The Diameter Replacement in Aldous--Fill Open Problem 4.49(ii) Fails}
\author{}
\date{16 September 2026}

\begin{document}
\maketitle

\begin{abstract}
Aldous and Fill record a multicommodity-flow estimate of the form
$\inf_{\mathbf f}\psi(\mathbf f)\le K\tau_2\log n$ and ask in Open Problem 4.49(ii) whether graph diameter can replace $\log n$ for a universal constant.  We give a negative construction.  Cones over sufficiently large fixed-degree bipartite Ramanujan graphs have diameter two and bounded relaxation time, while an explicit feasible metric in the weak dual forces the optimal stationary-demand congestion to grow without bound.  More precisely, for every $M>0$ there is a finite connected simple unweighted graph $G$ with
\[
 \frac{F(G)}{\tau_{\rm rel}(G)D(G)}>M,
\]
where $F(G)=\inf_{\mathbf f}\psi(\mathbf f)$ in the Aldous--Fill normalization.
\end{abstract}

\section{Source normalization and statement}
Let $G=(V,E)$ be a finite connected simple unweighted graph, let $P$ be simple random walk, and let
\[
 \pi_x=\frac{\deg(x)}{2|E|}.
\]
For every $a\in V$, let $f^{(a)}$ be an antisymmetric edge flow with divergence
\[
 \operatorname{div}f^{(a)}=\delta_a-\pi.
\]
Aldous--Fill define
\[
 \psi(\mathbf f)=
 \max_{\{i,j\}\in E}
 \frac{\sum_a\pi_a|f_{ij}^{(a)}|}{\pi_iP_{ij}},
 \qquad
 F(G):=\inf_{\mathbf f}\psi(\mathbf f).
\]
Write
\[
 \tau_{\rm rel}(G)=\frac1{1-\lambda_2(P)},
 \qquad
 D(G)=\operatorname{diam}(G),
\]
where $\lambda_2(P)$ is the second largest eigenvalue, not the second largest eigenvalue in absolute value.

Aldous--Fill equation (4.49) has the form
\[
 F(G)\le K\tau_{\rm rel}(G)\log |V|,
\]
and Open Problem 4.49(ii) asks whether $\log |V|$ can be replaced by $D(G)$ for a universal finite $K$.

\begin{theorem}\label{thm:main}
For every $M>0$ there exists a finite connected simple unweighted graph $G$ such that
\[
 \frac{F(G)}{\tau_{\rm rel}(G)D(G)}>M.
\]
Consequently no universal constant $K$ can make the diameter replacement in Aldous--Fill Open Problem 4.49(ii) valid for all such graphs.
\end{theorem}

\section{A weak metric dual bound}
For an unoriented edge $e=\{u,v\}$ put
\[
 c_e:=\pi_uP_{uv}=\pi_vP_{vu}.
\]
For nonnegative edge lengths $\ell=(\ell_e)$, let $d_\ell$ denote the induced shortest-path metric.

\begin{lemma}[Weak metric duality]\label{lem:dual}
If
\[
 \sum_{e\in E}c_e\ell_e\le1,
\]
then
\[
 F(G)\ge \sum_{a,b}\pi_a\pi_b d_\ell(a,b).
\]
\end{lemma}

\begin{proof}
Fix $a$ and set $\phi_a(v)=-d_\ell(a,v)$.  On every edge $e=\{u,v\}$,
$|\phi_a(u)-\phi_a(v)|\le\ell_e$.  Summation by parts against
$\operatorname{div}f^{(a)}=\delta_a-\pi$ gives
\[
 \sum_b\pi_b d_\ell(a,b)
 \le \sum_e \ell_e |f_e^{(a)}|.
\]
Multiply by $\pi_a$ and sum over $a$.  By the definition of $\psi$,
\[
 \sum_a\pi_a|f_e^{(a)}|\le \psi(\mathbf f)c_e.
\]
Therefore
\[
 \sum_{a,b}\pi_a\pi_b d_\ell(a,b)
 \le \psi(\mathbf f)\sum_ec_e\ell_e
 \le \psi(\mathbf f).
\]
Taking the infimum over feasible flow families proves the claim.
\end{proof}

\section{Ramanujan bases and cone geometry}
Fix an integer $s\ge4$ and set
\[
 d=s^2,\qquad r=s.
\]
Marcus--Spielman--Srivastava proved that for every degree $d>2$ there are infinite families of $d$-regular bipartite Ramanujan graphs.  Choose a connected simple member $H$ with $n$ vertices so large that
\[
 n\ge2B(d,r),
 \qquad
 B(d,r):=1+d\sum_{j=0}^{r-2}(d-1)^j.
\]
The selected 2-lift family may be taken connected: a disconnected 2-lift of a connected $d$-regular graph would produce a new adjacency eigenvalue $d$, whereas the selected new eigenvalues satisfy $|\theta|\le2\sqrt{d-1}<d$.  Simplicity is preserved by a 2-lift of a simple seed.

For independent uniform vertices $U,V$ of $H$, the ball bound implies
\[
 \Pr\{d_H(U,V)\ge r\}\ge\frac12,
\]
so
\begin{equation}\label{eq:truncdist}
 \E\min\{d_H(U,V),r\}\ge\frac r2.
\end{equation}

Form the cone
\[
 G=K_1\vee H
\]
with hub $h$.  Then
\[
 |E(G)|=\frac{n(d+2)}2,
 \qquad
 D(G)=2.
\]
The hub has degree $n$, every base vertex has degree $d+1$, and hence
\begin{equation}\label{eq:pi}
 \pi_h=\frac1{d+2},
 \qquad
 \pi_v=\frac{d+1}{n(d+2)}\quad(v\in H).
\end{equation}
Every physical edge has conductance
\begin{equation}\label{eq:cond}
 c_e=\frac1{n(d+2)}.
\end{equation}

\section{Relaxation time of the cone}
Let $A_Hv=\theta v$ with $v\perp\mathbf 1$ on the base.  Extending $v$ by zero at the hub gives a cone transition eigenvector with eigenvalue
\[
 \frac{\theta}{d+1}.
\]
On the two-dimensional subspace constant on the base together with the hub coordinate, the remaining eigenvalues are
\[
 1,\qquad -\frac1{d+1}.
\]
Since $H$ is Ramanujan,
\[
 |\theta|\le2\sqrt{d-1}
\]
for every nontrivial adjacency eigenvalue.  Therefore
\[
 \lambda_2(P_G)\le\frac{2\sqrt{d-1}}{d+1}
 <\frac{2s}{s^2+1},
\]
and hence
\begin{equation}\label{eq:gap}
 \frac1{\tau_{\rm rel}(G)}=1-\lambda_2(P_G)
 >\frac{(s-1)^2}{s^2+1}.
\end{equation}

\section{An exact feasible metric and the congestion lower bound}
Assign every spoke the length
\[
 x:=\frac{r(d+2)}{d+r}
\]
and every base edge the length
\[
 y:=\frac{2(d+2)}{d+r}.
\]
Then
\[
 2x=ry
\]
and, using \eqref{eq:cond},
\[
 \sum_ec_e\ell_e
 =\frac{nx+(nd/2)y}{n(d+2)}=1.
\]
Thus Lemma~\ref{lem:dual} applies.

The induced metric satisfies
\[
 d_\ell(h,v)=x,
 \qquad
 d_\ell(u,v)=y\min\{d_H(u,v),r\}
 \quad(u,v\in H).
\]
Indeed, any path using the hub pays at least two spokes, of total cost $2x=ry$, while every hub-free path lies in $H$ and costs at least $y d_H(u,v)$.

Let
\[
 \alpha:=\pi(H)=\frac{d+1}{d+2}.
\]
Using both ordered hub--base directions and \eqref{eq:truncdist}, Lemma~\ref{lem:dual} yields
\begin{align*}
 F(G)
 &\ge 2\frac1{d+2}\alpha x
 +\alpha^2y\E\min\{d_H(U,V),r\}\\
 &\ge
 \frac{r(d+1)(d+3)}{(d+2)(d+r)}.
\end{align*}
Combining this with $D(G)=2$ and \eqref{eq:gap}, then substituting $d=s^2$ and $r=s$, gives
\begin{equation}\label{eq:Ls}
 \frac{F(G)}{\tau_{\rm rel}(G)D(G)}
 >L_s:=
 \frac{(s-1)^2(s^2+3)}{2(s+1)(s^2+2)}.
\end{equation}
Finally,
\begin{equation}\label{eq:diverge}
 L_s-\frac{s-3}{2}
 =\frac{5s^2-2s+9}{2(s+1)(s^2+2)}>0.
\end{equation}
Hence the ratio in Theorem~\ref{thm:main} is larger than $(s-3)/2$ and is therefore unbounded.  This proves Theorem~\ref{thm:main}.

\section{Exact arithmetic certificate and audit status}
The accompanying script `check.py` uses Python's exact `fractions.Fraction` arithmetic to verify the length-budget identities, the simplification leading to \eqref{eq:Ls}, and the identity \eqref{eq:diverge} at several values of $s$.  In particular,
\[
 L_{64}=\frac{5422977}{177580}=30.5382193940759\ldots .
\]
The finite arithmetic is not a substitute for the proof or the external existence theorem.

The repository contains two independent claim-scoped audits of the theorem used here:
\begin{itemize}
\item `P26-AUDIT-C31-AF449-DIRECT-DUALITY-02 = PASS`;
\item `C30-AUDIT-C31-AF449-RAMANUJAN-CONE-01 = PASS`.
\end{itemize}
The present document is a packaging of that audited source-level claim.  No novelty or priority claim is made; a bounded literature search on 16 September 2026 found no prior published resolution of the exact Open Problem 4.49(ii), but no exhaustive bibliographic database audit was available.

\begin{thebibliography}{9}
\bibitem{AF}
D. Aldous and J. Fill,
\emph{Reversible Markov Chains and Random Walks on Graphs},
Chapter 4, Section 4.3.5 and Open Problem 4.49(ii),
\url{https://www.stat.berkeley.edu/~aldous/RWG/Book_Ralph/Ch4.S7.html}.

\bibitem{MSS}
A. W. Marcus, D. A. Spielman, and N. Srivastava,
Interlacing families I: Bipartite Ramanujan graphs of all degrees,
\emph{Annals of Mathematics} 182 (2015), 307--325.
\end{thebibliography}

\end{document}
